\documentclass[12pt,a4paper,oneside]{article}
\usepackage[brazil]{babel}
\usepackage[utf8]{inputenc}
\usepackage{olymp}

\setcounter{problem}{2}

\begin{document}

\begin{problem}{Elevador}{elevador}{2}
 
A \textit{United Floors Viçosa} (UFV) é uma empresa tradicional, com mais de 50 anos de experiência na fabricação de elevadores. Todos os projetos da UFV seguem as
mais estritas normas de segurança, mas infelizmente uma série de acidentes com seus elevadores manchou a reputação da empresa.

Ao estudar os acidentes, os engenheiros da companhia concluíram que, em vários casos, o acidente foi causado pelo excesso de passageiros no elevador. Por isso, a
UFV decidiu fiscalizar com mais rigor o uso de seus elevadores: foi instalado um sensor em cada porta que detecta a quantidade de pessoas que saem e entram
em cada andar do elevador.

A UFV tem os registros do sensor de todo um dia de funcionamento do elevador (que sempre começa vazio). Ela sabe que as pessoas são educadas e sempre deixam todos os passageiros que irão sair em um andar saírem antes de outros passageiros entrarem no elevador. Mas ainda assim ela tem tido dificuldade em decidir se a capacidade máxima do elevador foi excedida ou não.

%\Notes
%
%Observações, se tiver.

\InputFile

A primeira linha da entrada contém dois inteiros $N$ e $C$, indicando o número de leituras realizadas pelo sensor e a capacidade máxima do elevador, respectivamente ($1 \leq N \leq 1000$ e $1 \leq C \leq 1000$). As $N$ linhas seguintes contêm, cada uma, uma leitura do sensor. Cada uma dessas linhas contém dois inteiros $S$ e $E$, indicando quantas pessoas saíram e quantas pessoas entraram naquele andar, respectivamente ($0 \leq S \leq 1000$ e $0 \leq E \leq 1000$).


\OutputFile

Seu programa deve gerar apenas uma linha na saída, contendo o caractere `\texttt{S}', caso a capacidade do elevador tenha sido excedida em algum momento, ou o caractere `\texttt{N}', caso contrário.

\Example

\begin{example}
\exmp{
5 10
0 5
2 7
3 3
5 2
7 0
}{
N
}%
\end{example}

\begin{example}
\exmp{
5 10
0 3
0 5
0 2
3 4
6 4
}{
S
}%
\end{example}

\begin{example}
\exmp{
6 4
0 5
3 5
4 5
1 0
1 1
1 1
}{
S
}%
\end{example}

%Comentários sobre o exemplo, se necessário.
\footnotesize Problema da OBI 2010 - Olimpíada Brasileira de Informática

\end{problem}

\end{document}
